\documentclass[aspectratio=43,sanserif,professionalfonts]{beamer}
% valid values for aspectratio are 43 (4:3), 169 (16:9), 1610 (16:10)

% Default/core packages
\usepackage{graphicx}
\usepackage[english]{babel} % or [german]
\usepackage{helvet}
\usepackage[T1]{fontenc}
\usepackage{tikz}
\usepackage{calc}
\usepackage{hyperref}
\usepackage{xifthen}
\usepackage{etoolbox}
\usepackage{adjustbox}

\usepackage{amsmath}
\usepackage{mathtools}
\usepackage{algorithm2e}
\usepackage{tcolorbox}

\colorlet{mygreen}{green!50!black}

\DeclareMathOperator{\rot}{rot}
\DeclareMathOperator*{\argmin}{arg\,min}

% ====================

% Apply i8 styling
% ============
% Configure Beamer Theme

% Some libs
\usetikzlibrary{positioning}
\usetikzlibrary{shapes.misc}

% Disable navigation symbols
\beamertemplatenavigationsymbolsempty

% Primary RWTH colors

\definecolor{rwth-blue}{RGB}{0, 84, 159}
\definecolor{rwth-black}{RGB}{0, 0, 0}
\definecolor{rwth-magenta}{RGB}{227, 0, 102}
\definecolor{rwth-yellow}{RGB}{255, 237, 0}
\definecolor{rwth-petrol}{RGB}{0, 97, 101}
\definecolor{rwth-turquoise}{RGB}{0, 152, 161}
\definecolor{rwth-turquois}{RGB}{0, 152, 161}
\definecolor{rwth-green}{RGB}{87, 171, 39}
\definecolor{rwth-may-green}{RGB}{189, 205, 0}
\definecolor{rwth-orange}{RGB}{246, 168, 0}
\definecolor{rwth-red}{RGB}{204, 7, 30}
\definecolor{rwth-bordeaux}{RGB}{161, 16, 53}
\definecolor{rwth-violet}{RGB}{97, 33, 88}
\definecolor{rwth-purple}{RGB}{122, 111, 172}

% Faded RWTH colors

\definecolor{rwth-blue-100}{RGB}{0, 84, 159}
\definecolor{rwth-blue-75}{RGB}{64, 127, 183}
\definecolor{rwth-blue-50}{RGB}{142, 186, 229}
\definecolor{rwth-blue-25}{RGB}{199, 221, 242}
\definecolor{rwth-blue-10}{RGB}{232, 241, 250}

\definecolor{rwth-black-100}{RGB}{0, 0, 0}
\definecolor{rwth-black-75}{RGB}{100, 101, 103}
\definecolor{rwth-black-50}{RGB}{156, 158, 159}
\definecolor{rwth-black-25}{RGB}{207, 209, 210}
\definecolor{rwth-black-10}{RGB}{236, 237, 237}

\definecolor{rwth-magenta-100}{RGB}{227, 0, 102}
\definecolor{rwth-magenta-75}{RGB}{233, 96, 136}
\definecolor{rwth-magenta-50}{RGB}{214, 158, 177}
\definecolor{rwth-magenta-25}{RGB}{249, 210, 218}
\definecolor{rwth-magenta-10}{RGB}{253, 238, 240}

\definecolor{rwth-yellow-100}{RGB}{255, 237, 0}
\definecolor{rwth-yellow-75}{RGB}{255, 240, 85}
\definecolor{rwth-yellow-50}{RGB}{255, 245, 155}
\definecolor{rwth-yellow-25}{RGB}{255, 250, 209}
\definecolor{rwth-yellow-10}{RGB}{255, 253, 238}

\definecolor{rwth-petrol-100}{RGB}{0, 97, 101}
\definecolor{rwth-petrol-75}{RGB}{45, 127, 131}
\definecolor{rwth-petrol-50}{RGB}{125, 164, 167}
\definecolor{rwth-petrol-25}{RGB}{191, 208, 209}
\definecolor{rwth-petrol-10}{RGB}{230, 236, 236}

\definecolor{rwth-turquoise-100}{RGB}{0, 152, 161}
\definecolor{rwth-turquoise-75}{RGB}{0, 177, 183}
\definecolor{rwth-turquoise-50}{RGB}{137, 204, 207}
\definecolor{rwth-turquoise-25}{RGB}{202, 231, 231}
\definecolor{rwth-turquoise-10}{RGB}{235, 246, 246}

\definecolor{rwth-turquois-100}{RGB}{0, 152, 161}
\definecolor{rwth-turquois-75}{RGB}{0, 177, 183}
\definecolor{rwth-turquois-50}{RGB}{137, 204, 207}
\definecolor{rwth-turquois-25}{RGB}{202, 231, 231}
\definecolor{rwth-turquois-10}{RGB}{235, 246, 246}

\definecolor{rwth-green-100}{RGB}{87, 171, 39}
\definecolor{rwth-green-75}{RGB}{141, 192, 96}
\definecolor{rwth-green-50}{RGB}{184, 214, 152}
\definecolor{rwth-green-25}{RGB}{221, 235, 206}
\definecolor{rwth-green-10}{RGB}{242, 247, 236}

\definecolor{rwth-may-green-100}{RGB}{189, 205, 0}
\definecolor{rwth-may-green-75}{RGB}{208, 217, 92}
\definecolor{rwth-may-green-50}{RGB}{224, 230, 154}
\definecolor{rwth-may-green-25}{RGB}{240, 243, 208}
\definecolor{rwth-may-green-10}{RGB}{249, 250, 237}

\definecolor{rwth-orange-100}{RGB}{246, 168, 0}
\definecolor{rwth-orange-75}{RGB}{250, 190, 80}
\definecolor{rwth-orange-50}{RGB}{253, 212, 143}
\definecolor{rwth-orange-25}{RGB}{254, 234, 201}
\definecolor{rwth-orange-10}{RGB}{255, 247, 234}

\definecolor{rwth-red-100}{RGB}{204, 7, 30}
\definecolor{rwth-red-75}{RGB}{216, 92, 65}
\definecolor{rwth-red-50}{RGB}{230, 150, 121}
\definecolor{rwth-red-25}{RGB}{243, 205, 187}
\definecolor{rwth-red-10}{RGB}{250, 235, 227}

\definecolor{rwth-bordeaux-100}{RGB}{161, 16, 53}
\definecolor{rwth-bordeaux-75}{RGB}{182, 82, 86}
\definecolor{rwth-bordeaux-50}{RGB}{205, 139, 135}
\definecolor{rwth-bordeaux-25}{RGB}{229, 197, 192}
\definecolor{rwth-bordeaux-10}{RGB}{245, 232, 229}

\definecolor{rwth-violet-100}{RGB}{97, 33, 88}
\definecolor{rwth-violet-75}{RGB}{131, 78, 117}
\definecolor{rwth-violet-50}{RGB}{168, 133, 158}
\definecolor{rwth-violet-25}{RGB}{210, 192, 205}
\definecolor{rwth-violet-10}{RGB}{237, 229, 234}

\definecolor{rwth-purple-100}{RGB}{122, 111, 172}
\definecolor{rwth-purple-75}{RGB}{155, 145, 193}
\definecolor{rwth-purple-50}{RGB}{188, 181, 215}
\definecolor{rwth-purple-25}{RGB}{222, 218, 235}
\definecolor{rwth-purple-10}{RGB}{242, 240, 247}

% Specific VCI design colors

\definecolor{i8-titlepage-top}{RGB}{21, 77, 156}
\definecolor{i8-titlepage-bottom}{RGB}{17, 60, 120}

\definecolor{i8-title-top}{RGB}{0, 109, 199}
\definecolor{i8-title-bottom}{RGB}{0, 80, 145}

\definecolor{i8-nr-top}{RGB}{0, 109, 199}
\definecolor{i8-nr-bottom}{RGB}{0, 83, 150}

% Templates

% no outer theme
%\useoutertheme{infolines}
% Inner Theme
\useinnertheme{rectangles}

% ============
% Fonts
\usefonttheme[onlymath]{serif}
% \renewcommand{\familydefault}{\sfdefault}

% ============
% Colors

\setbeamercolor{title}{fg=rwth-blue}
\setbeamercolor{titlelike}{fg=rwth-blue}

\setbeamercolor{normal text}{fg=rwth-black}

\setbeamercolor{alerted text}{fg=rwth-red}

\setbeamercolor{section in toc}{fg=rwth-blue}
\setbeamercolor{subsection in toc}{fg=rwth-black}

\setbeamercolor{structure}{fg=rwth-blue}

% ============
% Helper

\newrobustcmd*{\footfullcitenomark}{%
    \AtNextCite{%
        \let\thefootnote\relax
        \let\mkbibfootnote\mkbibfootnotetext}%
    \footfullcite}

% ============
% Size

\fboxsep0pt

\newcommand{\rwthMarginLeftRight}{0.038\paperwidth}
\newcommand{\rwthMarginRightLogo}{0.022\paperwidth}
\newcommand{\rwthMarginFooterText}{0.0385\paperheight}
\newcommand{\rwthFooterTextLeft}{0.10\paperwidth}
\newcommand{\rwthFooterLogoWidth}{0.35\paperwidth}
\newcommand{\rwthFooterLogoHeight}{0.153\paperheight}
\newcommand{\footerHeight}{0.16\paperheight}
\newlength{\headerHeight}
\newlength{\frameHeight}

\setbeamersize{text margin left=\rwthMarginLeftRight,text margin right=\rwthMarginLeftRight}

% ============
% Header
\newcommand{\rwthInsertFrameTitle}{

    %\color{black} \hrulefill
    \begin{tikzpicture}[overlay]
    \color{white}
    \draw (0,0) circle [radius=0pt];
    \shade[rounded corners=5pt, top color=i8-title-top, bottom color=i8-title-bottom] (-.25,10pt) rectangle  +(\paperwidth,-1);
    % \color{black}
    % \draw[line width=0.5pt] (0,0) -- (\paperwidth-\rwthMarginLeftRight-\rwthMarginLeftRight,0);
    \end{tikzpicture}
    
    \vspace{-25pt}
    \parbox[top][1.2cm][b]{\textwidth}{
        \color{white}
        \textbf{\Large \insertframetitle}
        \vphantom{p} % force "p-line"
    }
    \vspace{2pt}

    \ifx\insertframesubtitle\empty
    \else
    \vspace{0.6em}
    \color{black}
    \textbf{\large \insertframesubtitle}
    \vspace{0.5em}
    \fi
}
\setbeamertemplate{frametitle}{
    \rwthInsertFrameTitle
}

\newcommand{\updateFrameHeight}
{
    \settototalheight\headerHeight{\vbox{\rwthInsertFrameTitle}}
    \addtolength{\headerHeight}{1.6em}
    \setlength{\frameHeight}{\textheight-\headerHeight}
}

% patch frame(sub)title to automatically update \frameHeight
% BREAKS cites, multi-frame frames :(
%\let\beamerOldFrametitle\frametitle
%\renewcommand{\frametitle}[1]{\beamerOldFrametitle{#1}\updateFrameHeight}
%\let\beamerOldFramesubtitle\framesubtitle
%\renewcommand{\framesubtitle}[1]{\beamerOldFramesubtitle{#1}\updateFrameHeight}

% ============
% Footer

% Extra data
\newcommand{\footlineextra}[1]{\gdef\insertfootlineextra{#1}}
\footlineextra{}

% Remove footnote rule
\renewcommand*{\footnoterule}{}

\def\rwthOptionTitlePage{}
\def\rwthOptionNoNumbering{}
\setbeamertemplate{footline}
{
    \begin{tikzpicture}[overlay]    
        % show rest only if non-lean
        \ifthenelse{\equal{\rwthOptionTitlePage}{lean}}{}
        {
            % page nr
            \shade[rounded corners=5pt, top color=i8-nr-top, bottom color=i8-nr-bottom] (\rwthMarginLeftRight-.25cm,-\rwthMarginFooterText) rectangle +(1, +1.7);
        }
    \end{tikzpicture}
    \begin{tikzpicture}
        % anchor coord system
        \color{white}
        \draw (0,0) circle [radius=0pt];

        % logo
        \node[below left](Footerlogo) at (\paperwidth-\rwthMarginRightLogo, 0) {\includegraphics[height=\rwthFooterLogoHeight]{i8style/rwth_vci_rgb_white.png}};

        % show rest only if non-lean
        \ifthenelse{\equal{\rwthOptionTitlePage}{lean}}{}
        {
            % page nr
            \color{white}
            \tiny
            \ifx\rwthOptionNoNumbering\empty
                \node[below] at (\rwthMarginLeftRight+.18cm,-\rwthMarginFooterText-3pt) {\textbf{\centering\large\insertframenumber}};
            \fi

            % footer text
            \color{rwth-blue}
            \tiny
            \node[below right, text width=\paperwidth-\rwthMarginLeftRight-\rwthFooterLogoWidth-\rwthFooterTextLeft] at (\rwthFooterTextLeft,-\rwthMarginFooterText) {
                \insertshortauthor\\
                \insertshorttitle\\
                \insertdate
                };

            % horizontal line
            % \color{black}
            % \draw[line width=0.5pt] (\rwthMarginLeftRight,0) -- (\paperwidth-\rwthMarginLeftRight,0);
        }

    \end{tikzpicture}
}

% Custom Title Pages

\newcommand{\rwthSubtitleOrAuthor}{\ifx\insertsubtitle\empty\insertauthor\else\insertsubtitle\fi}

\newcommand{\rwthTitlepageColor}[1][lean]{
    \def\rwthOptionTitlePage{#1}
    \begin{frame}[noframenumbering]
        \begin{tikzpicture}[remember picture,overlay]
            % colored tile
            \node[inner sep=0pt] (topleft) at (current page.north west) {};
            \node[below = 0.3333\paperheight-2pt of topleft,inner sep=0pt] (tileend) {};
            \shade[top color=i8-titlepage-top,bottom color=i8-titlepage-bottom,inner sep=0pt] (tileend) rectangle +(\paperwidth, 0.3333\paperheight);

            % Title
            \node[below right = 0.3333\paperheight+0.0167\paperheight and 1em of topleft] (titlestart) {
                \color{rwth-blue}
                \adjustbox{max width=\textwidth}{\textbf{\Large \inserttitle}}
                };
            % Subtitle
            \node[below right] (subtitlestart) at (titlestart.south west) {
                \color{black}
                \large \rwthSubtitleOrAuthor
                };
        \end{tikzpicture}
    \end{frame}
    \def\rwthOptionTitlePage{}
}

\newcommand{\rwthTitlepagePicture}[2][lean]{
    \def\rwthOptionTitlePage{#1}
    \begin{frame}[noframenumbering]
        \begin{tikzpicture}[remember picture,overlay]
        % colored tile
        \node[inner sep=0pt] (topleft) at (current page.north west) {};
        \node[below = 0.3333\paperheight-2pt of topleft,inner sep=0pt] (tileend) {};
        %\fill[rwth-blue] (tileend) rectangle +(\paperwidth, 0.3333\paperheight);

        % Title img
        \node[above right,inner sep=0pt] at (tileend) {\includegraphics[width=\paperwidth]{#2}};

        % Title
        \node[below right = 0.3333\paperheight+0.0167\paperheight and 1em of topleft] (titlestart) {
            \color{rwth-blue}
            \adjustbox{max width=\textwidth}{\textbf{\Large \inserttitle}}
        };
        % Subtitle
        \node[below right] (subtitlestart) at (titlestart.south west) {
            \color{black}
            \large \rwthSubtitleOrAuthor
        };
        \end{tikzpicture}
    \end{frame}
    \def\rwthOptionTitlePage{}
}

\newcommand{\rwthTitlepagePictureLarge}[2][lean]{
    \def\rwthOptionTitlePage{#1}
    \begin{frame}[noframenumbering]
        \begin{tikzpicture}[remember picture,overlay]
        % colored tile
        \node[inner sep=0pt] (topleft) at (current page.north west) {};
        \node[below = 0.6666\paperheight-2pt of topleft,inner sep=0pt] (tileend) {};
        %\fill[rwth-blue] (tileend) rectangle +(\paperwidth, 0.6666\paperheight);

        % Title img
        \node[above right,inner sep=0pt] at (tileend) {\includegraphics[width=\paperwidth]{#2}};

        % Title
        \node[below right = 0.6666\paperheight+0.0167\paperheight and 13pt of topleft] (titlestart) {
            \color{rwth-blue}
            \adjustbox{max width=\textwidth}{\textbf{\Large \inserttitle}}
        };
        % Subtitle
        \node[below right] (subtitlestart) at (titlestart.south west) {
            \color{black}
            \large \rwthSubtitleOrAuthor
        };
        \end{tikzpicture}
    \end{frame}
    \def\rwthOptionTitlePage{}
}


% ====================
% Title
\title{Instant Field Aligned Meshes}

% ====================
% Author
\author{Til Mohr}
\institute[Visual Computing Group]
{
    Computer Graphics and Multimedia \\
    Prof.~Dr.~Leif~Kobbelt \\
    RWTH Aachen University
}

% ====================
% Date
\date{February 9, 2022}

% ====================
% Subject
\subject{Computer Graphics}

% ====================
% Extra footer lines
\footlineextra{}

% ====================
% Start document
\begin{document}
\rwthTitlepagePictureLarge{img/title.png} % banner image 2/3 of slide size
% alternative title pages:
%\rwthTitlepagePicture{img/title.png} % banner image 1/3 of slide size
% \rwthTitlepageColor % blue background 1/3 of slide size

\begin{frame}
	\frametitle{Remeshing?}
	\begin{columns}
		\begin{column}{0.5\textwidth}
			\center
			\includegraphics[width=\textwidth]{img/input.png}
		\end{column}
		\begin{column}{0.5\textwidth}
			\center
			\includegraphics[width=\textwidth]{img/output-pure.png}
		\end{column}
	\end{columns}
\end{frame}

\begin{frame}
\frametitle{Triangle Meshes}
	\center
	\includegraphics[width=\textwidth]{img/triangles.png}
\end{frame}

\begin{frame}
	\frametitle{Quad Meshes}
	\begin{columns}
		\begin{column}{0.5\textwidth}
			\centering
			Quad-dominant Mesh
		\end{column}
		\begin{column}{0.5\textwidth}
			\centering
			Pure Quad Mesh
		\end{column}
	\end{columns}
	\includegraphics[width=\textwidth]{img/Quad.png}
\end{frame}

\begin{frame}
	\frametitle{Uniformity \& Isotropy}
	\begin{tcolorbox}[colback=green!5,colframe=green!40!black,title=Uniformity]
		Faces are of the same size
	\end{tcolorbox}
	\pause
	\begin{tcolorbox}[colback=green!5,colframe=green!40!black,title=Isotropy]
		Uniformity of directions
	\end{tcolorbox}
	\pause
	$\Rightarrow$ Uniform and isotropic mesh:
	\begin{itemize}
		\item Equilateral triangles (Triangle Mesh)
		\item Squares (Quad Mesh)
	\end{itemize}
\end{frame}

\begin{frame}
	\frametitle{Remeshing}
	\begin{tcolorbox}[colback=green!5,colframe=green!40!black,title=Correlation Problem]
		Finding the location of a new vertex on the input surface
	\end{tcolorbox}
	\pause
	\centering
	~\\
	\begin{columns}
		\begin{column}{0.5\textwidth}
			\textbf{Global Remeshing:}\\
			Create a new output mesh\\
			~\\
			{\color{mygreen} $\Rightarrow$ High quality meshes and globally optimal}\\
			{\color{red} $\Rightarrow$ High time and design complexity}
		\end{column}
		\begin{column}{0.5\textwidth}
			\textbf{Local Remeshing:}\\
			Modify the input mesh connectivity\\
			~\\
			{\color{red} $\Rightarrow$ Lower quality meshes and not globally optimal}\\
			{\color{mygreen} $\Rightarrow$ Low time and design complexity}
			\\~\\
		\end{column}
	\end{columns}
\end{frame}

\iffalse
\begin{frame}
	\frametitle{Global Remeshing}
	\begin{tcolorbox}[colback=green!5,colframe=green!40!black,title=Mesh Parameterization]
		Mapping between meshes of the same surface
	\end{tcolorbox}
	\pause
	~\\
	Solves the correspondence problem\\
	{\color{mygreen} $\Rightarrow$ High quality meshes and globally optimal}\\
	\pause
	~\\
	Challenging to compute\\
	{\color{red} $\Rightarrow$ High time and design complexity}
\end{frame}

\begin{frame}
	\frametitle{Local Remeshing}
	Modify the initial mesh connectivity using topological operations\\
	\pause
	~\\
	Limited to local information around each vertex\\
	{\color{red} $\Rightarrow$ Lower quality meshes and not globally optimal}\\
	\pause
	~\\
	Easy to compute and parallelize \\
	{\color{mygreen} $\Rightarrow$ Low time and design complexity}
\end{frame}
\fi

\iffalse
\begin{frame}
	\frametitle{$N$-RoSy Fields}
	\only<-2,4->{
		\begin{tcolorbox}[colback=green!5,colframe=green!40!black,title=$N$-way Rotational Symmetry]
			Shape that is invariant under rotations by a multiple of $\frac{2\pi}{N}$
		\end{tcolorbox}
		\pause
		~\\
		Can be thought of as vector fields with $N$ vectors at each point\\
		\pause
		~\\
		Symmetries on surfaces:
		\begin{itemize}
			\item	\textbf{Intrinsic:} Preserve geodesic distances
			\item	\textbf{Extrinsic:} Preserve euclidean distances
		\end{itemize}

	}
	\only<3>{
		\centering
		\includegraphics[width=\textwidth]{img/n-Rosy-Singularity.png}
	}
\end{frame}
\fi

% METHOD

\begin{frame}
	\frametitle{Method}
	Combines ideas of global and local remeshing\\
	\pause
	~\\
	\textbf{Input:} mesh, point cloud, etc.\\
	~\\
	\pause
	\begin{enumerate}
		\item	Transform input into a graph
		\item	Calculate the orientation field
		\item	Calculate the position field
		\item	Export the mesh
	\end{enumerate}
	~\\
	\pause
	\textbf{Output:} uniform, isotropic triangle / quad(-based) mesh
\end{frame}

\begin{frame}
	\frametitle{Input Representation}
	Weighted directed graph $\mathcal{G} = (\mathcal{V}, \mathcal{E})$:
	\begin{itemize}
		\item	Vertex $i \in \mathcal{V}$: position and normal direction $v_i, n_i \in \mathbb{R}^3$
		\pause
		\item	$\mathcal{E} \subseteq \mathcal{V} \times \mathcal{V}$:
				\begin{itemize}
					\item	Meshes: Usual set of edges
					\item	Point Clouds: $K$-nearest neighbors by distance
				\end{itemize}
		\pause
		\item	$(i,j) \in \mathcal{E}$: weights $w_{ij}$ $\rightarrow$ affect orientation- \& position field
	\end{itemize}
	\pause
	~\\
	$\Rightarrow$ Simplicity allows for broad application
\end{frame}

\begin{frame}
	\frametitle{Fields}
	\begin{columns}
		\begin{column}{0.5\textwidth}
			\textbf{1. Orientation Field}\\
			\begin{itemize}
				\item	$s_o$-RoSy field
				\item	Guides edges in the output mesh
			\end{itemize}
		\end{column}
		\pause
		\begin{column}{0.5\textwidth}
			\textbf{2. Position Field}\\
			\begin{itemize}
				\item	Uses the orientation field
				\item	Specifies the location of vertices in the output mesh
			\end{itemize}
		\end{column}
	\end{columns}
	\pause
	~\\~\\
	Input parameters:
	\begin{itemize}
		\item	Type of mesh:
				\begin{itemize}
					\item	Triangle mesh: $s_o=6$
					\item	Quad-dominant mesh: $s_o=4$
				\end{itemize}
		\item	Desired edge length: $\rho$
		\item	Intrinsic or extrinsic optimization
	\end{itemize}
\end{frame}

\begin{frame}
	\frametitle{Orientation Field}
	\centering
	\includegraphics[width=\textwidth]{img/orientation-field-illustration.png}
\end{frame}

\begin{frame}
	\frametitle{Orientation Field}
	Idea:
	\begin{itemize}
		\item	Output mesh should be uniform and isotropic
		\item	Therefore there are rotational symmetries in the output mesh
		\item	Thus we compute an $s_o$-RoSy field
	\end{itemize}
	\pause
	~\\
	\begin{columns}
		\begin{column}{0.7\textwidth}
			Let $o_i \in \mathbb{R}^3$ be a representational direction vector for each $i \in \mathcal{V}$.\\

			\pause
			\begin{tcolorbox}[colback=green!5,colframe=green!40!black,title=Symmetry Group - Integer Rotations]
				$\mathcal{R}_{s_o}(o,n,k) \coloneqq \rot \left(n, k \cdot \frac{2\pi}{s_o} \right) o, \quad k \in \mathbb{Z}$ \\
				$\mathcal{R}_{s_o}(o,n) \coloneqq \mathcal{R}_{s_o}(o,n,0, \dots, s_o-1)$
			\end{tcolorbox}
		\end{column}
		\begin{column}{0.2\textwidth}
			\centering
			\includegraphics[width=\textwidth]{img/integer-rotation.png}
		\end{column}
	\end{columns}
\end{frame}

\begin{frame}
	\frametitle{Orientation Field}
	Orientation field $O$ consists of $\mathcal{R}_{s_o}(o_i,n_i)$ for every $i \in \mathcal{V}$\\
	$\Rightarrow$ Smooth out the orientation field\\
	$\Rightarrow$ Minimize angles between adjacent vertices\\
	\pause
	~\\
	Let $k_{ij} \in \{0, \dots, s_o-1\}$ for each $(i,j) \in \mathcal{E}$.\\
\end{frame}

\begin{frame}
	\frametitle{Orientation Field - Optimization}
	\includegraphics[width=\textwidth]{img/orientation-field-intrinsic-extrinsic-horizontal.png}

	~\\~\\
	\begin{columns}
		\begin{column}{0.5\textwidth}
			\textbf{Intrinsic Smoothness Energy:}\\
			$\sum\limits_{(i,j) \in \mathcal{E}} \angle(o_i, \mathcal{R}_{s_o}(o_{ji}, n_i, k_{ij}))^2$\\~\\~\\
		\end{column}
		\begin{column}{0.5\textwidth}
			\textbf{Extrinsic Smoothness Energy:}\\
			$\sum\limits_{(i,j) \in \mathcal{E}} \angle(\mathcal{R}_{s_o}(o_{i}, n_i, k_{ij}),$\\
			\qquad\qquad $\mathcal{R}_{s_o}(o_{j}, n_j, k_{ji}))^2$
		\end{column}
	\end{columns}
\end{frame}

\begin{frame}
	\frametitle{Orientation Field - Intrinsic Optimization}
	\begin{algorithm}[H]
		\For{$i \in \mathcal{V}$}
		{
			\For{$j \in \mathcal{N}(i)$}
			{
				\tcp{brute force "smoothest" $k_{ij}$}
				$k_{ij} \leftarrow \argmin\limits_{0 \leq k' < s_0} \angle(o_i, \mathcal{R}_{s_o}(o_{ji}, n_i, k'))$
			}
			\tcp{optimize $o_i$ by rotation}
			$o_i \leftarrow \sum\limits_{j \in \mathcal{N}(i)} w_{ij} \cdot \mathcal{R}_{s_o}(o_{ji}, n_i, k_{ij})$\\
			$o_i \leftarrow \frac{o_i}{\Vert o_i \Vert}$
		}
	\end{algorithm}
\end{frame}

\begin{frame}
	\frametitle{Orientation Field}
	\begin{columns}
		\begin{column}{0.5\textwidth}
			\center
			\includegraphics[width=\textwidth]{img/orientation.png}
		\end{column}
		\begin{column}{0.5\textwidth}
			\center
			\includegraphics[width=\textwidth]{img/output-pure.png}
		\end{column}
	\end{columns}
\end{frame}

\begin{frame}
	\frametitle{Position Field}
	\centering
	\includegraphics[width=\textwidth]{img/position-field-illustration.png}
\end{frame}

\begin{frame}
	\frametitle{Position Field}
	Idea:
	\begin{itemize}
		\item	Use the orientation field $O$ to determine the position of vertices in the output mesh
		\item	Per-vertex local parameterization with fixed gradient $o_i$
		\item	Space vertices along edges $\rho$ units apart
	\end{itemize}
	\pause
	~\\
	\begin{columns}
		\begin{column}{0.7\textwidth}
			\begin{tcolorbox}[colback=green!5,colframe=green!40!black,title=Symmetry Group - Integer Translations]
				$\mathcal{T}_{s_o}(p,n,o,t) \coloneqq p + \rho \cdot \sum\limits_{k=0}^{\frac{s_o}{2}-1} t_k \mathcal{R}_{s_o}(o,n,k)$ \\
				$\mathcal{T}_{s_o}(p,n,o) \coloneqq \left\{\mathcal{T}_{s_o}(p,n,o,t) \vert t \in \mathbb{Z}^{\frac{s_o}{2}}\right\}$
			\end{tcolorbox}
		\end{column}
		\begin{column}{0.2\textwidth}
			\centering
			\includegraphics[width=\textwidth]{img/integer-translation.png}
		\end{column}
	\end{columns}
\end{frame}

\begin{frame}
	\frametitle{Position Field}
	Position field $P$ consists of local parameterizations $\mathcal{T}_{s_o}(p_i,n_i,o_i)$ for every $i \in \mathcal{V}$\\
	$\Rightarrow$ Smooth out the position field\\
	$\Rightarrow$ Minimize distances between local parameterizations of adjacent vertices\\
	\pause
	~\\
	Let $t_{ij} \in \mathbb{Z}^{\frac{s_o}{2}}$ for each $(i,j) \in \mathcal{E}$.\\
\end{frame}

\begin{frame}
	\frametitle{Position Field - Optimization}
	\includegraphics[width=\textwidth]{img/position-field-intrinsic-extrinsic-horizontal.png}

	~\\~\\
	\begin{columns}
		\begin{column}{0.5\textwidth}
			\textbf{Intrinsic Smoothness Energy:}\\
			$\sum\limits_{(i,j) \in \mathcal{E}} \Vert p_i - \mathcal{T}_{s_o}(p_{ji}, n_i, o_i, t_{ij}) \Vert^2_2$\\~\\~\\
		\end{column}
		\begin{column}{0.5\textwidth}
			\textbf{Extrinsic Smoothness Energy:}\\
			$\sum\limits_{(i,j) \in \mathcal{E}} \Vert \mathcal{T}_{s_o}(p_i, n_i, o_i, t_{ij}) -$\\
			\qquad\quad $\mathcal{T}_{s_o}(p_j, n_j, o_j, t_{ji}) \Vert^2_2$
		\end{column}
	\end{columns}
\end{frame}

\begin{frame}
	\frametitle{Position Field - Intrinsic Optimization}
	\begin{algorithm}[H]
		\For{$i \in \mathcal{V}$}
		{
			\For{$j \in \mathcal{N}(i)$}
			{
				\tcp{brute force "smoothest" $t_{ij}$}
				$t_{ij} \leftarrow \argmin\limits_{t' \in Q_{ij}} \Vert p_i - \mathcal{T}_{s_o}(p_{ji}, n_i, o_i, t')) \Vert$
			}
			\tcp{optimize $p_i$ by translation}
			$p_i \leftarrow \frac{\sum\limits_{j \in \mathcal{N}(i)} w_{ij} \cdot \mathcal{T}_{s_o}(p_{ji}, n_i, o_i, t_{ij})}{\sum\limits_{j \in \mathcal{N}(i)} w_{ij}}$\\
		}
	\end{algorithm}
\end{frame}

\begin{frame}
	\frametitle{Position Field}
	\begin{columns}
		\begin{column}{0.5\textwidth}
			\center
			\includegraphics[width=\textwidth]{img/position.png}
		\end{column}
		\begin{column}{0.5\textwidth}
			\center
			\includegraphics[width=\textwidth]{img/output-pure.png}
		\end{column}
	\end{columns}
\end{frame}

\begin{frame}
	\frametitle{Mesh Extraction}
	Idea:
	\begin{enumerate}
		\item	Create new undirected graph $\mathcal{G}' \coloneqq (\mathcal{V}', \mathcal{E}')$
		\item	Add all $p_i$ to $\mathcal{V'}$
		\pause
		\item	Connect $(p_i, p_j)$ if $t_{ij}$ is unit vector-valued\\
				$\Rightarrow$ $p_i$ and $p_j$ are just apart by $\rho$
		\pause
		\item	Additional filtering to add edges
	\end{enumerate}
	\pause
	~\\
	$\mathcal{G}'$ has now the final structure of the mesh
\end{frame}

% Comparison
\begin{frame}
	\frametitle{Comparison - Quality}
	\centering
	\includegraphics[width=0.6\textwidth]{img/quality.png}
\end{frame}

\begin{frame}
	\frametitle{Comparison - Computation Time}
	\centering
	\includegraphics[width=\textwidth]{img/computation-time.png}
\end{frame}

\begin{frame}
	\frametitle{Conclusion}
	\begin{itemize}
		\item	Quality issues (singularities) need to be dealt with
		\item	Linear computation time makes it great for practical usage
		\item	Simplicity in design and modularity\\
				$\Rightarrow$ Optimization steps can be used in other work
	\end{itemize}
\end{frame}

% References
\begin{frame}[allowframebreaks]
	\frametitle{References}
	\nocite{*}
	\bibliographystyle{amsalpha}
	\bibliography{references}
\end{frame}

\end{document}
